% !TEX program = lualatex
\documentclass[aspectratio=169,professionalfonts]{beamer}
\newcommand{\slidemode}{1}  % カラー投影モード
\usepackage{beamer_template}
\usepackage{enumerate}

\newcommand{\LectureNum}{14}
\newcommand{\LectureTitle}{スペクトル}
\newcommand{\CourseName}{応用数学}
\renewcommand{\TermName}{前期}

\begin{document}

\begin{frame}{第14回　スペクトル}
  \begin{block}{本回の内容}
  スペクトル
  \end{block}
  \vfill\hfill{\scriptsize \textcolor{gray}{第3章 フーリエ解析　§2}}
\end{frame}

\begin{frame}{例2　（p.\,100）}
  \begin{exampleblock}{問題}
    周期4の矩形波のスペクトルを求めよ\\[2pt]
    $f(x) = \begin{cases} 1 & (-1 < x < 1) \\ 1/2 & (x = 1, 3) \\ 0 & (1 < x < 3) \end{cases},  f(x+4) = f(x)$\\[2pt]
  \end{exampleblock}
\end{frame}

\begin{frame}{例2　解答}
  \begin{exampleblock}{解答}
  \end{exampleblock}
  {\footnotesize \textcolor{gray}{\textit{線スペクトル（ $\omega$ の値がとびとびに現れる）}}}
\end{frame}

\begin{frame}{例3　（p.\,101）}
  \begin{exampleblock}{問題}
    次の関数のスペクトルを求めよ\\[2pt]
    $f(x) = \begin{cases} 1 & (|x| < 1) \\ 1/2 & (|x| = 1) \\ 0 & (|x| > 1) \end{cases}$\\[2pt]
  \end{exampleblock}
\end{frame}

\begin{frame}{例3　解答}
  \begin{exampleblock}{解答}
  \end{exampleblock}
  {\footnotesize \textcolor{gray}{\textit{連続スペクトル（ $S_i(\omega)$ の値が連続的に変化する）}}}
\end{frame}

\begin{frame}{演習問題（p.\,100--103）}
  \begin{exampleblock}{自分で解いてみよう}
  \begin{description}
    \item[問10] 次の関数のスペクトルを求めよ
    $f(x) = 1 - |x|  (|x| \leq 1),  f(x+2) = f(x)$\\
    \item[問11] 次の関数のスペクトルを求めよ
    $f(x) = { 1-|x|  (|x| \leq 1),  0  (|x| > 1) }$\\
  \end{description}
  \end{exampleblock}
\end{frame}

\begin{frame}{練習問題2　（p.\,104）}
  \begin{block}{}
  \begin{description}
    \item[問1] 次の関数のフーリエ変換を求めよ
    $(1) f(x) = { 2  (0 \leq x < 2),  1  (2 \leq x < 3),  0  (x < 0, x \geq 3) }$\\
    $(2) g(x) = { 1 - x/2  (0 \leq x < 2),  0  (x < 0, x \geq 2) }$\\
    \item[問2] 次の関数について、(1) はフーリエ正弦変換、(2) はフーリエ余弦変換をそれぞれ求めよ
    $(1) f(x) = { \sin x  (|x| \leq \pi),  0  (|x| > \pi) }$\\
    $(2) g(x) = { \cos x  (|x| \leq \pi/2),  0  (|x| > \pi/2) }$\\
  \end{description}
  \end{block}
\end{frame}

\end{document}
